\documentclass[../../../include/open-logic-section]{subfiles}

\begin{document}

\olfileid{sth}{cardinals}{cp}
\olsection{مبدأ كانتور}

ارجع إلى~\olref[ordinals][vn]{sec}. فقد كنا نبحث المجموعات المرتبة
ترتيبًا جيدًا، فرأينا أن نجعل لها أشياء تقوم مقام تلك الترتيبات. ومن
هنا أدخلنا الأعداد الترتيبية، ثم أثبتنا في
\olref[ordinals][ordtype]{ordtypesworklikeyouwant} أنها تؤدي المقصود:
\[
	\ordtype{A, <} = \ordtype{B, \lessdot}
	\text{ إذا وفقط إذا كان } \tuple{A, <} \isomorphic \tuple{B, \lessdot}.
\]
ثم ارجع إلى ما هو أسبق، إلى~\olref[sfr][siz][equ]{sec}. هنالك، ونحن
نعمل بسذاجة، أدخلنا معنى «حجم» المجموعة؛ فقلنا إن مجموعتين متساويتان
عدديًا،~$\cardeq{A}{B}$، إذا وفقط إذا وجدت !!a{bijection}
$f \colon A \to B$. وهذا المعنى {أبسط} في جوهره من معنى الترتيب
الجيد؛ لأن العبرة هنا بوجود !!{bijection}s لا بكون تلك الـ!!{bijection}s
حافظة لبنية ما،
كما يشترط في تماثلات الترتيب.

ويهدينا ذلك إلى أمر ظاهر: فكما جعلنا \emph{الأعداد الترتيبية} موازين
للترتيبات الجيدة، نجعل أشياء مخصوصة، هي \emph{الأعداد الكاردينالية}،
موازين للحجم. وهذا هو غرض الفصل.

وقبل تعريف هذه الأعداد ينبغي أن نحدد الأصل الذي نطلب منها تحقيقه.
إذا رمزنا بـ$\card{X}$ إلى كاردينالية~$X$، أردنا:
\[
	\card{A} = \card{B} \text{ إذا وفقط إذا كان } \cardeq{A}{B}.
\]
ونسمي هذا \emph{مبدأ كانتور}؛ فإن كانتور، فيما يظهر، أول من تصوره
تصورًا تام الوضوح. وسيأتي مزيد كلام في صلته \emph{بمبدأ هيوم} في
\olref[hp]{sec}. فالمطلوب إذن أن نعرّف~$\card{X}$ لكل~$X$ على وجه
يحقق مبدأ كانتور.

\end{document}
